% =====================================================================
\chapter{Rectas, planos, sistemas lineales y determinantes}
\label{tema:2}
% =====================================================================

Una ecuación lineal en tres incógnitas describe un plano, y resolver un sistema de tres ecuaciones es cortar tres planos. Este tema escribe rectas y planos con vectores, resuelve sistemas pequeños por eliminación gaussiana e interpreta el determinante como un área o un volumen con signo. Todo se limita a tamaños $2\times2$ y $3\times3$: la teoría general (rango, independencia lineal, determinantes de orden $n$) es materia de los módulos~2 y~3 del curso. La presentación sigue las sesiones~5--16 de la unidad~1 de 18.02SC~\cite{mit1802}, la §2.5 de OpenStax~\cite{openstax3}, las secciones~11.2--11.4 de Strang~\cite{strangcalc} y, para la eliminación, el capítulo~2 de su \emph{Introduction to Linear Algebra}~\cite{strangla}.

\section{Rectas y planos}

\begin{definicion}[Recta]\label{def:recta}
La \emph{recta} que pasa por el punto $P_0$ con vector director $\vv\neq\vcero$ es el conjunto de puntos
\[
\vr(t)=P_0+t\,\vv,\qquad t\in\R .
\]
En coordenadas, $x=x_0+tv_1$, $y=y_0+tv_2$, $z=z_0+tv_3$ (\emph{ecuaciones paramétricas}). La recta que pasa por $P$ y $Q$ tiene vector director $\overrightarrow{PQ}$.
\end{definicion}

En el plano $\R^2$ una recta admite además una ecuación \emph{implícita} $ax+by=c$; el vector $(a,b)$ es perpendicular a la recta, porque si $(x_1,y_1)$ y $(x_2,y_2)$ son dos de sus puntos, restando las ecuaciones queda $(a,b)\cdot(x_2-x_1,y_2-y_1)=0$. En el espacio la ecuación implícita de una sola ecuación describe un plano.

\begin{definicion}[Plano]\label{def:plano}
El \emph{plano} que pasa por $P_0=(x_0,y_0,z_0)$ con \emph{vector normal} $\vn=(a,b,c)\neq\vcero$ es el conjunto de los puntos $X=(x,y,z)$ tales que $\vn\cdot(X-P_0)=0$, es decir,
\[
a(x-x_0)+b(y-y_0)+c(z-z_0)=0\quad\Longleftrightarrow\quad ax+by+cz=d,\qquad d=\vn\cdot P_0 .
\]
\end{definicion}

Recíprocamente, toda ecuación $ax+by+cz=d$ con $(a,b,c)\neq\vcero$ es un plano de normal $(a,b,c)$: basta tomar uno de sus puntos como $P_0$. Para el plano que pasa por tres puntos no alineados $P$, $Q$, $R$ se toma como normal $\vn=\overrightarrow{PQ}\times\overrightarrow{PR}$, que es perpendicular a ambos vectores (\cref{prop:vectorial}).

\addtocounter{definicion}{1}%

\begin{proposicion}[Posición relativa de dos planos]\label{prop:dos-planos}
Sean $\Pi_1:\vn_1\cdot X=d_1$ y $\Pi_2:\vn_2\cdot X=d_2$.
\begin{enumerate}
  \item Si $\vn_1$ y $\vn_2$ son proporcionales, los planos son paralelos: coinciden si las ecuaciones son proporcionales y no tienen puntos comunes en caso contrario.
  \item Si no lo son, se cortan en una recta de vector director $\vn_1\times\vn_2$.
\end{enumerate}
El \emph{ángulo} entre $\Pi_1$ y $\Pi_2$ es el ángulo agudo (o recto) $\varphi$ entre sus normales: $\cos\varphi=\abs{\vn_1\cdot\vn_2}/(\norm{\vn_1}\norm{\vn_2})$.
\end{proposicion}

\addtocounter{definicion}{1}%

Para cortar una recta $\vr(t)=P_0+t\vv$ con un plano $\vn\cdot X=d$ se sustituye $\vr(t)$ en la ecuación del plano y se despeja $t$. Por ejemplo, la recta $(1+t,\,t,\,t)$ y el plano $x+2y+3z=7$ dan $1+6t=7$, $t=1$, y el punto de corte es $(2,1,1)$. Si $\vn\cdot\vv=0$, la recta es paralela al plano: o no lo corta o está contenida en él.

\section{Sistemas lineales y eliminación gaussiana}\label{sec:t2-gauss}

\begin{definicion}[Sistema lineal y matriz]\label{def:sistema}
Un sistema de $m$ ecuaciones lineales con $n$ incógnitas ($m,n\le3$ en este tema) se escribe $A\vx=\vb$, donde $A=(a_{ij})$ es la \emph{matriz de coeficientes}, $\vx$ el vector de incógnitas y $\vb$ el de términos independientes. El producto $A\vx$ es el vector cuya componente $i$ es el producto escalar de la fila $i$ de $A$ por $\vx$. El sistema es \emph{compatible determinado} si tiene una única solución, \emph{compatible indeterminado} si tiene infinitas e \emph{incompatible} si no tiene ninguna.
\end{definicion}

Más en general, el producto de matrices $AB$ (definido cuando el número de columnas de $A$ coincide con el de filas de $B$) tiene como elemento $(i,j)$ el producto escalar de la fila $i$ de $A$ por la columna $j$ de $B$. La matriz identidad $I$ tiene unos en la diagonal y ceros fuera, y cumple $AI=IA=A$.

\begin{definicion}[Operaciones elementales]\label{def:elementales}
Las \emph{operaciones elementales} sobre las ecuaciones de un sistema son: (E1) intercambiar dos ecuaciones; (E2) multiplicar una ecuación por un número no nulo; (E3) sumar a una ecuación un múltiplo de otra.
\end{definicion}

\begin{proposicion}\label{prop:elementales}
Las operaciones elementales no cambian el conjunto de soluciones de un sistema.
\end{proposicion}

El \emph{método de Gauss} aplica operaciones elementales hasta llegar a un sistema \emph{escalonado} (cada ecuación tiene al menos una incógnita inicial menos que la anterior) y lo resuelve por \emph{sustitución regresiva}, de abajo arriba. En la práctica se trabaja con la \emph{matriz ampliada} $(A\mid\vb)$.

\addtocounter{definicion}{1}%

Si al escalonar aparece una fila $(0\ 0\ 0\mid c)$ con $c\neq0$, el sistema es incompatible; si aparece una fila entera de ceros y quedan menos ecuaciones que incógnitas, las incógnitas que no encabezan ninguna fila son \emph{libres} y el sistema es compatible indeterminado. Geométricamente, en un sistema $3\times3$ cada ecuación es un plano: la solución única corresponde a tres planos con un único punto común; la indeterminación, a planos que comparten una recta (o un plano); la incompatibilidad, a planos sin punto común a los tres.

\section{Determinantes de orden 2 y 3}

\begin{definicion}[Determinantes]\label{def:det3}
\[
\det\begin{pmatrix}a&b\\c&d\end{pmatrix}=\begin{vmatrix}a&b\\c&d\end{vmatrix}=ad-bc,
\qquad
\begin{vmatrix}a_1&a_2&a_3\\b_1&b_2&b_3\\c_1&c_2&c_3\end{vmatrix}
=a_1\begin{vmatrix}b_2&b_3\\c_2&c_3\end{vmatrix}-a_2\begin{vmatrix}b_1&b_3\\c_1&c_3\end{vmatrix}+a_3\begin{vmatrix}b_1&b_2\\c_1&c_2\end{vmatrix}.
\]
\end{definicion}

Comparando con la \cref{def:vectorial}, el determinante $3\times3$ de filas $\va,\vb,\mathbf c$ es exactamente el producto mixto $\va\cdot(\vb\times\mathbf c)$.

\begin{teorema}[Interpretación geométrica]\label{teo:det-area}
\begin{enumerate}
  \item $\abs{ad-bc}$ es el área del paralelogramo de lados $(a,b)$ y $(c,d)$ en $\R^2$.
  \item El valor absoluto del determinante $3\times3$ de filas $\va,\vb,\mathbf c$ es el volumen del paralelepípedo de aristas $\va,\vb,\mathbf c$.
\end{enumerate}
\end{teorema}

\begin{proposicion}[Propiedades de los determinantes]\label{prop:det-propiedades}
Para matrices $2\times2$ y $3\times3$:
\begin{enumerate}
  \item intercambiar dos filas cambia el signo del determinante; si dos filas son iguales, el determinante es nulo;
  \item multiplicar una fila por $\lambda$ multiplica el determinante por $\lambda$;
  \item sumar a una fila un múltiplo de otra no cambia el determinante;
  \item $\det(A^{\mathsf T})=\det A$ y $\det(AB)=\det A\cdot\det B$.
\end{enumerate}
\end{proposicion}

Para orden~2 todas las propiedades se comprueban con la fórmula (véase el \cref{eje:2-producto-det}); para orden~3 se deducen de las propiedades del producto mixto, y la demostración general, para cualquier orden, se da en el módulo~3 del curso. La propiedad~(iii) dice que la eliminación gaussiana con operaciones (E3) conserva el determinante: el de una matriz escalonada es el producto de los elementos de su diagonal. En el \cref{ej:t2-gauss}, $\det A=1\cdot1\cdot2=2$.

\begin{teorema}[Sistemas $2\times2$ y regla de Cramer]\label{teo:cramer}
El sistema $ax+by=e$, $cx+dy=f$ tiene solución única si y solo si $\Delta=ad-bc\neq0$, y en ese caso
\[
x=\frac{1}{\Delta}\begin{vmatrix}e&b\\f&d\end{vmatrix},\qquad y=\frac{1}{\Delta}\begin{vmatrix}a&e\\c&f\end{vmatrix}.
\]
\end{teorema}

Lo mismo vale en orden~3: un sistema $3\times3$ tiene solución única si y solo si $\det A\neq0$, es decir, si las tres normales (las filas de $A$) no son coplanarias (\cref{prop:volumen}); la regla de Cramer sustituye la columna $j$ de $A$ por $\vb$ para obtener la incógnita $j$.

\begin{definicion}[Matriz inversa]\label{def:inversa}
Una matriz cuadrada $A$ es \emph{invertible} si existe $A^{-1}$ con $AA^{-1}=A^{-1}A=I$. Entonces el sistema $A\vx=\vb$ tiene la única solución $\vx=A^{-1}\vb$.
\end{definicion}

Para orden~2, si $ad-bc\neq0$,
\begin{equation}\label{eq:t2-inversa}
\begin{pmatrix}a&b\\c&d\end{pmatrix}^{-1}=\frac{1}{ad-bc}\begin{pmatrix}d&-b\\-c&a\end{pmatrix},
\end{equation}
como se comprueba multiplicando. Una matriz es invertible si y solo si su determinante es no nulo.

\addtocounter{definicion}{1}%

\addtocounter{definicion}{1}%

